% Level 2, Week 3 -- Lists: Collections of Things
\lesson{3}{Lists: Collections of Things}{A list keeps many values in order, and the first one is number 0.}
\lessoninfo{Concepts: lists, index, append, len, for x in list (I), random.choice (R)}{Tools: Thonny, word cards}{CSTA: 2-AP-11, 3A-AP-14, 2-IC-21}

\begin{objectives}
  \item The student can make a list, add to it with \texttt{append} and count it with \texttt{len}. (L2.G4; L2.W03.1)
  \item The student can get an item by its index, starting from 0, and explain why \texttt{items[len(items)]} gives an \texttt{IndexError}. (L2.G4; L2.W03.2)
  \item The student can loop over a list with \texttt{for item in items} and with \texttt{for i in range(len(items))}. (L2.G2; L2.W03.3)
  \item The student can choose a color-blind-safe palette and label a picture so it never relies on color alone. (L2.G9; L2.W03.4)
\end{objectives}

\begin{materials}
  \item Print (one per student): workbook chapter 3, homework 3.
  \item \enquote{Human list} cards (3 minutes): five cards reading \enquote{lantern}, \enquote{rope}, \enquote{map}, \enquote{apple}, \enquote{key},
        and five sticky notes numbered 0, 1, 2, 3, 4 for the floor. Fewer than 4 students? Stand cards up on chairs or hold them with soft toys.
  \item On each Mac: open \texttt{2-intermediate/exercises/week03/ex-palette-stripes/starter.py} in Thonny (\keys{Cmd}+\keys{O}).
  \item A timer, paper and pencils, and optionally colored pencils for the Wrap-up.
\end{materials}

\needspace{16\baselineskip}
\section{At a glance}
Times are a guide: set a timer for each phase and move on when it rings.
\begin{timeplan}
  \timeslot{0--8}{Warm-up}{Unplugged \enquote{Human list}: index from 0, \texttt{append}, \texttt{len} and a live \texttt{IndexError}.}
  \timeslot{8--18}{Concept}{Lists in Python: square brackets, index, \texttt{append}, \texttt{len}, two kinds of \texttt{for} loop. Designing for everyone: color-blind-safe palettes.}
  \timeslot{18--28}{Guided demo}{Palette Stripes: append two colors, then turn one stripe into a loop that draws them all.}
  \timeslot{28--52}{Hands-on build}{Pairs: Exercise 3.1 Palette Stripes (finish), 3.2 Adventure Inventory, then 3.3 Top 3 High Scores. Swap every 8 minutes. Last 5 minutes: free-build.}
  \timeslot{52--60}{Wrap-up}{Show the stripes, accessibility check, exit ticket, reflection, homework 3.}
\end{timeplan}

\section{Warm-up (0--8 min)}
\begin{unplugged}[Human list]
Stick the numbered notes on the floor in a row: 0, 1, 2, 3. Four students (or chairs) hold \enquote{lantern}, \enquote{rope}, \enquote{map}, \enquote{apple}
and stand on them. Say: \enquote{This line is a list called \texttt{backpack}.} Then call out, and let them act:
\begin{itemize}[leftmargin=1.2em]
  \item \enquote{\texttt{backpack[0]}, step forward!} (lantern) \quad \enquote{\texttt{backpack[2]}!} (map)
  \item \enquote{\texttt{len(backpack)}?} Everyone counts: 4.
  \item \enquote{\texttt{backpack.append("key")}}: a new card joins at the \emph{end}, on note 4.
  \item \enquote{\texttt{backpack[5]}!} Nobody is there. Everyone shouts \enquote{IndexError!}
\end{itemize}
Ask: \enquote{With 5 things in the list, what is the biggest index?} (4: always one less than the length, because we start at 0.)
\end{unplugged}

\section{Concept (8--18 min)}
\begin{teacherscript}[Say]
\enquote{So far each variable held one value. A \textbf{list} holds many values in order, inside square brackets, separated by commas.
Each item has a position number called its \textbf{index}, and Python starts counting at 0, just like our floor.
\texttt{append} adds to the end, \texttt{len} counts the items, and a \texttt{for} loop visits every item in turn.}
\end{teacherscript}

\begin{pycode}[Lists in four lines]
backpack = ["lantern", "rope", "map"]
backpack.append("apple")            # ["lantern", "rope", "map", "apple"]
print(len(backpack), backpack[0], backpack[-1])   # 4 lantern apple
for item in backpack:               # visits every item, in order
    print(item)
\end{pycode}

\begin{analogy}[Lockers in a hallway]
A list is a row of lockers. The first locker has the number 0 painted on it. To get something you need the locker number (the index),
not the thing's name. If the hallway has 4 lockers, they are numbered 0 to 3: ask for locker 4 and the custodian says
\enquote{that locker does not exist}. That is Python's \texttt{IndexError}. \texttt{-1} is a shortcut for \enquote{the last locker}.
\end{analogy}

\begin{safety}[Color-blind-safe palettes: designing for everyone]
Say: \enquote{About 1 in 12 boys and 1 in 200 girls see some colors differently, often red and green. If a game shows danger
only in red and safety only in green, those players cannot tell them apart.} Show the Okabe-Ito palette in the workbook: these
8 colors stay different for almost everyone. Rule for our art and games: use these colors \textbf{and} add a label, number,
shape or pattern, so color is never the only clue. Our stripes are labeled with their index number for exactly this reason.
\end{safety}

\section{Guided demo (18--28 min)}
\begin{workedexample}[From one stripe to a loop of stripes]
\begin{enumerate}[leftmargin=1.4em]
  \item Open \texttt{ex-palette-stripes/starter.py} and press \menu{Run} (\keys{F5}). One orange stripe labeled 0 appears.
  \item TODO 1: under the first list, type \texttt{palette.append("\#0072B2")} and \texttt{palette.append("\#D55E00")}. Run: still one stripe.
        Ask: \enquote{The list has 5 colors now. Why only one stripe?} (Because \texttt{draw\_stripes()} only draws \texttt{colors[0]}.)
  \item TODO 3: replace the last line of \texttt{draw\_stripes()} with a loop. Type it slowly and say each part aloud:
\begin{pycode}
    for i in range(len(colors)):
        draw_stripe(t, x, 100, width, 200, colors[i], str(i))
        x = x + width
\end{pycode}
        \enquote{\texttt{len(colors)} is 5, so \texttt{i} goes 0, 1, 2, 3, 4. Each time we draw stripe \texttt{i} and move right.}
        Run: five labeled stripes.
  \item Ask: \enquote{Why \texttt{str(i)}?} (\texttt{write} shows text; \texttt{str} turns the number into text, like the types from week 1.)
  \item Leave TODO 4 (print the first and last color) for the pairs.
\end{enumerate}
\end{workedexample}
The finished Palette Stripes, for you:
\pyfile{../exercises/week03/ex-palette-stripes/solution.py}

\section{Hands-on build (28--52 min)}
\begin{pairroles}[Driver, navigator and the index checker]
The \textbf{driver} types; the \textbf{navigator} keeps a paper sketch of the list with index numbers under each item and
checks every \texttt{[ ]}. Swap at the 8-minute timer.
\textbf{1 student:} you are the navigator with the paper sketch. \textbf{3 students:} add a \textbf{tester} who runs
\texttt{check\_my\_work.py} after each TODO. \textbf{4 students:} two pairs; pair A starts at 3.2 and pair B at 3.3, then swap and explain.
\end{pairroles}

\begin{enumerate}[leftmargin=1.4em]
  \item \textbf{Exercise 3.1 Palette Stripes}: finish TODO 4 (and TODO 2 for a sixth color). About 4 minutes.
  \item \textbf{Exercise 3.2 Adventure Inventory} (\texttt{week03/ex-inventory}, about 10 minutes). Starts with this week's bug hunt.
  \item \textbf{Exercise 3.3 Top 3 High Scores} (\texttt{week03/ex-high-scores}, about 6 minutes, or for fast finishers).
        New tricks: \texttt{sorted(scores, reverse=True)} and \texttt{best[:3]}.
\end{enumerate}

\begin{bughunt}[IndexError: one past the end]
Adventure Inventory crashes with \texttt{IndexError: list index out of range} in \texttt{last\_item()}, which uses
\texttt{backpack[len(backpack)]}. Do not give the answer. Ask: \enquote{How many items are in the backpack? What are their index
numbers? So which index is the last one?} Let the navigator draw the four boxes 0--3. Two good fixes: \texttt{backpack[len(backpack) - 1]}
or \texttt{backpack[-1]}. This is called an \textbf{off-by-one} bug: one of the most common bugs in all of programming.
\end{bughunt}

\begin{freebuild}[Pattern studio]
Last 5 minutes: make your own stripes. Ideas: thinner stripes with \texttt{width=30}; the palette twice with \texttt{palette * 2};
your own order of colors; a rainbow \enquote{flag} for your hero team. Keep labels on every stripe.
\end{freebuild}

\section{Wrap-up (52--60 min)}
Show the stripes. Do a quick accessibility check together: \enquote{If you could not see color at all, could you still tell the stripes apart?}
(Yes: the numbers.)

\begin{formativecheck}[Exit ticket]
\begin{pycode}
pets = ["cat", "dog", "fish"]
pets.append("frog")
print(len(pets), pets[1], pets[-1])
\end{pycode}
What does it print? Which index would cause an \texttt{IndexError}?
\textbf{Answer:} \texttt{4 dog frog}. Any index from 4 up (the indexes are 0--3).
\textbf{Look for:} \enquote{cat} for \texttt{pets[1]} means the student still counts from 1. Repeat the floor game next week for 2 minutes.
\end{formativecheck}

\begin{reflection}[Say it or write it]
\begin{itemize}[leftmargin=1.2em]
  \item Why do you think Python starts counting at 0? Does the locker analogy help?
  \item What is one list you could use in a game you would like to make?
  \item How did we make our stripes work for people who see colors differently?
\end{itemize}
\end{reflection}
Hand out homework 3 (Loot Bag).

\section{Notes for the teacher}
\begin{commonmistakes}
  \item \textbf{Counting from 1.} \texttt{items[1]} is the \emph{second} item. The first is \texttt{items[0]}.
  \item \textbf{One past the end.} \texttt{items[len(items)]} always crashes. The last item is \texttt{items[len(items) - 1]} or \texttt{items[-1]}.
  \item \textbf{Forgetting the brackets or quotes.} \texttt{["lantern", rope]} gives \texttt{NameError}: text needs quotes.
  \item \textbf{\texttt{append} with an equals sign.} \texttt{backpack = backpack.append("map")} sets the list to \texttt{None}. Just write \texttt{backpack.append("map")}.
  \item \textbf{Mixing up the loops.} In \texttt{for item in backpack}, \texttt{item} is the thing itself; in \texttt{for i in range(len(backpack))}, \texttt{i} is the index number.
  \item \textbf{x not moving.} If every stripe lands in the same place, \texttt{x = x + width} is missing or not indented inside the loop.
\end{commonmistakes}

\begin{differentiation}
\supportpath Keep the floor numbers on the table during the build. In Palette Stripes, type the loop line together and let the
pair type the body. Skip Exercise 3.3; it can be a homework extra. In Inventory, give \texttt{has\_item} as
\texttt{return item in backpack} (a shortcut) so the pair can focus on the index bug.
\extensionpath In Inventory, write \texttt{drop\_item(backpack, item)} with \texttt{backpack.remove(item)} and protect it with
\texttt{has\_item}. In High Scores, store a made-up nickname with each score in a second list and print \texttt{1st: 300 (Comet)}.
\end{differentiation}

\begin{stretchchallenge}[Checkerboard palette]
Draw two rows of stripes, where the second row uses the colors in reverse order. Tip: \texttt{list(reversed(palette))} or
\texttt{palette[::-1]}. Can you label every square with its index in its own row?
\end{stretchchallenge}

\begin{teachernote}[If you are short on time]
Shorten the human list to \texttt{[0]}, \texttt{len} and the IndexError shout. Do 3.1 and 3.2 only; 3.3 becomes an optional homework extra.
Keep the color-blind-safe talk short but do not skip the accessibility check in the Wrap-up.
\end{teachernote}
